% SPDX-License-Identifier: CC-BY-4.0
% Copyright (c) 2026 Martin Schr\"oder <info@swedishembedded.com>

\section{Method}
\label{sec:method}

\subsection{When does a Fourier basis matter at all?}
\label{sec:absorption}

A decoder transformer keeps a residual stream $x_t\in\R^{d}$ per token. Every
map applied to that stream is either a learned dense linear map (the
attention projections $W_Q,W_K,W_V,W_O$ and the feed-forward projections) or
a normalization. Let $Q\in\R^{d\times d}$ be any orthogonal matrix, for
instance a real Fourier basis, and write the stream in that basis, $z_t=Q^\top
x_t$.

\begin{proposition}[Basis absorption]
\label{prop:absorption}
Let $f$ be a pre-norm decoder whose only operations on the residual stream are
(i) addition, (ii) dense linear maps $W$ into or out of the stream and (iii)
RMS normalization $\operatorname{RMSNorm}_\gamma(x)=\gamma\odot x/\operatorname{rms}(x)$.
Then $f$ with residual stream $z=Q^\top x$ computes the same function as $f$
with residual stream $x$ after the substitutions $W_{\mathrm{in}}\mapsto
W_{\mathrm{in}}Q$, $W_{\mathrm{out}}\mapsto Q^\top W_{\mathrm{out}}$, token
embedding $E\mapsto EQ$, and $\operatorname{diag}(\gamma)\mapsto
Q^\top\operatorname{diag}(\gamma)Q$ for every normalization gain.
\end{proposition}

\begin{proof}
$\operatorname{rms}(Qz)=\operatorname{rms}(z)$ because $Q$ is orthogonal, so
$\operatorname{RMSNorm}_\gamma(Qz)=Q\,(Q^\top\operatorname{diag}(\gamma)Q)\,z/\operatorname{rms}(z)$.
A map reading the stream sees $W_{\mathrm{in}}x=W_{\mathrm{in}}Qz$; a map
writing it adds $W_{\mathrm{out}}u$ to $x$, that is $Q^\top W_{\mathrm{out}}u$
to $z$. Addition commutes with $Q$. Induction over the layers gives the claim.
\end{proof}

Two consequences fix the design and the controls of this paper. First, a
model that merely stores its hidden state as Fourier coefficients and keeps
dense learned maps everywhere is a conventional transformer: the basis is
absorbed into the weights, apart from the normalization gain, which becomes a
dense symmetric matrix instead of a diagonal one. Second, keeping only $m<d$
coefficients of such a model is a conventional transformer of width $m$. A
``rotated'' transformer and a ``truncated-spectrum'' transformer with dense
maps therefore need no new architecture: they are a conventional decoder of
width $d$ and of width $m$, and those are the controls we train
(\cref{sec:protocol}). A Fourier basis can only change what a model computes
through an operator that is defined relative to the basis: a map that is
structured in it (diagonal per frequency mode) or a nonlinearity applied in
the signal domain rather than the coefficient domain. The spectral mixer below
has both.

\subsection{The truncated real Fourier representation}

Fix a signal length $D$ (even) and a retained width $m\le D$. The real
orthonormal Fourier basis of $\R^D$, ordered by frequency, is
\begin{equation}
\label{eq:basis}
\varphi_0(j)=\tfrac{1}{\sqrt D},\quad
\varphi_{2k-1}(j)=\sqrt{\tfrac2D}\cos\tfrac{2\pi kj}{D},\quad
\varphi_{2k}(j)=\sqrt{\tfrac2D}\sin\tfrac{2\pi kj}{D},\quad
\varphi_{D-1}(j)=\tfrac{(-1)^j}{\sqrt D},
\end{equation}
for $k=1,\dots,D/2-1$ and $j=0,\dots,D-1$. The synthesis matrix
$S=[\varphi_0,\dots,\varphi_{m-1}]\in\R^{D\times m}$ keeps the $m$ lowest
modes; its columns are orthonormal, $S^\top S=I_m$, and the analysis map is
$S^\top$. A coefficient vector $z\in\R^m$ represents the band-limited signal
$Sz\in\R^D$, and Parseval's identity $\lVert Sz\rVert_2=\lVert z\rVert_2$
holds exactly. We count capacity in real degrees of freedom: $m$ coefficients
are $m$ real numbers, with no conjugate-symmetric duplicate stored.

\subsection{Spectral language model}

The spectral language model (SpectralLM) is a pre-norm decoder whose residual
stream is $z_t\in\R^m$, the retained coefficients of a virtual signal of
length $D$. It never materializes a $D$-wide residual stream.

\paragraph{Embedding and head.} A table $E\in\R^{V\times m}$ maps tokens
directly to coefficients, and the head computes logits $E_{\mathrm{out}}
\operatorname{RMSNorm}(z_t)$ with $E_{\mathrm{out}}\in\R^{V\times m}$. No
transform is involved.

\paragraph{Attention.} Each block applies grouped-query causal attention with
rotary position embedding \citep{su2021roformer} and per-head QK
normalization, as in Qwen3 \citep{qwen3report}, to $\operatorname{RMSNorm}(z_t)$.
By \cref{prop:absorption} attention is indifferent to the basis, so it is
identical in SpectralLM and in the width-$m$ control, and its key-value cache
holds $m$-wide coefficient states by construction.

\paragraph{Spectral mixer.} The feed-forward sublayer is replaced by a mixer
with $C$ channels. For $h=\operatorname{RMSNorm}(z)\in\R^m$ and learned
per-mode filters $G,U,W\in\R^{C\times m}$,
\begin{equation}
\label{eq:mixer}
M(h)=\sum_{c=1}^{C} W_c\odot S^\top\!\Big[\silu\big(S(G_c\odot h)\big)\odot S(U_c\odot h)\Big].
\end{equation}
\begin{figure}[ht]
\centering
\begin{tikzpicture}[node distance=4mm and 7mm, every node/.style={font=\small},
  box/.style={draw, rounded corners=2pt, minimum height=7mm, inner sep=3pt, align=center},
  op/.style={draw, circle, inner sep=1pt}]
\node[box] (h) {$h$\\[-1pt]\scriptsize $[m]$};
\node[box, right=of h, yshift=6mm] (g) {$G_c\odot h$\\[-1pt]\scriptsize $[C,m]$};
\node[box, right=of h, yshift=-6mm] (u) {$U_c\odot h$\\[-1pt]\scriptsize $[C,m]$};
\node[box, right=of g] (sg) {$S(\cdot)$\\[-1pt]\scriptsize $[C,D]$};
\node[box, right=of u] (su) {$S(\cdot)$\\[-1pt]\scriptsize $[C,D]$};
\node[box, right=of sg, yshift=-6mm] (glu) {$\silu(a)\odot b$\\[-1pt]\scriptsize $[C,D]$};
\node[box, right=of glu] (an) {$S^\top(\cdot)$\\[-1pt]\scriptsize $[C,m]$};
\node[box, right=of an] (y) {$\sum_c W_c\odot t_c$\\[-1pt]\scriptsize $[m]$};
\draw[-{Stealth}] (h) -- (g); \draw[-{Stealth}] (h) -- (u);
\draw[-{Stealth}] (g) -- (sg); \draw[-{Stealth}] (u) -- (su);
\draw[-{Stealth}] (sg) -- node[above, font=\scriptsize] {$a$} (glu); \draw[-{Stealth}] (su) -- node[below, font=\scriptsize] {$b$} (glu);
\draw[-{Stealth}] (glu) -- (an); \draw[-{Stealth}] (an) -- node[above, font=\scriptsize] {$t$} (y);
\end{tikzpicture}
\caption{The spectral mixer for one row. Per-mode filters act on the $m$
retained coefficients, synthesis $S$ moves to the $D$-sample signal domain
where the gated product couples frequencies, analysis $S^\top$ returns to $m$
coefficients, and per-mode output weights sum the $C$ channels.}
\label{fig:mixer}
\end{figure}

Read from the inside out (\cref{fig:mixer}): each channel filters the input spectrum twice
(diagonal maps, $2m$ parameters), synthesizes both filtered spectra into
$D$-sample signals, applies the SwiGLU product \citep{shazeer2020glu} pointwise
in the signal domain, analyses the result back to $m$ coefficients, and
weights it per mode. The mixer has $3Cm$ parameters, against $3mf$ for a
SwiGLU of hidden width $f$.

The signal-domain product is where frequencies interact. Pointwise
multiplication of two signals is circular convolution of their spectra, so a
mode of the output depends on pairs of input modes whose frequencies add or
subtract to it; no diagonal map could produce that coupling, and it is what a
per-mode filter bank alone lacks. If both factors are band-limited to the
lowest $m$ modes their product occupies at most the lowest $2m$, so with
$D\ge 2m$ the bilinear part of \cref{eq:mixer} is computed without aliasing
before $S^\top$ projects it back. The $\silu$ factor is not band-limited, so
$D/m$ is a hyperparameter that trades aliasing against cost rather than a
condition for exactness.

\paragraph{Relation to Fourier neural operators.} \Cref{eq:mixer} is a
one-dimensional Fourier neural operator layer \citep{li2020fno} along the
embedding axis rather than along space or sequence: per-mode weights in the
spectral domain, a pointwise nonlinearity in the signal domain, and fixed
transforms between them. It differs in three respects: the operator runs on
the feature axis of each token independently, its spectral weights are
diagonal per channel instead of a per-mode channel-mixing matrix, and its
nonlinearity is a gated product rather than a pointwise activation of a
linear lift.

\subsection{Cost}
\label{sec:cost}

Per token and layer, ignoring normalization, a SwiGLU of hidden width $f$
costs $6mf$ floating-point operations and $3mf$ parameters. The mixer costs
$O(Cm)$ operations for its filters, $O(CD)$ for the signal-domain product, and
$3C$ transforms of length $D$: two syntheses and one analysis per channel. Computed as dense products with $S$,
the transforms cost $6CDm$ operations in total; computed with a fast real
Fourier transform they cost $O(CD\log D)$. The mixer therefore saves
parameters at any $C<f$, but saves arithmetic only with a fast transform: at
equal arithmetic, a dense-product mixer can afford only $C=f/D$ channels,
while an FFT mixer can afford $C=\Theta(fm/(D\log D))$. This is why the
implementation (\cref{sec:implementation}) carries both transform paths and
gates the fast one against the dense one.
